\section{Directed transition bounds for compressed overlap search}
\label{sec:lcs-transitions}

\subsection{Equal alphabets can still have different local order}

The same-alphabet control of Section~\ref{sec:lcs-bounds},
$(ab)^Nac,(ab)^Nca$, exhausted the work allowance at $N=2^{500}$.
An alphabet-only bound cannot distinguish these words, but their adjacent
letter pairs differ. Any common substring must use pairs that occur in
both words. The new stage applies this stronger necessary condition after
the alphabet-run stage has failed to finish the query.

It is also necessary to improve the preceding whole-donor stage. A fast
partial-overlap query is of little use if an impossible full-donor query
has already consumed its budget. We therefore use directed pair sets to
reject impossible whole-donor orientations before building inverse words
or invoking the occurrence matcher. Both uses are exact pruning; the
existing complete string algorithms remain available when the conditions
are insufficient.

\subsection{Exact directed pairs from grammar cuts}

For a word $W$, let $E(W)$ be its set of ordered adjacent signed-letter
pairs. This is a linear-word definition; no cyclic closing edge is added
unless explicitly requested by a caller. For each reachable binary rule
$U=AB$, record the pair
$(\operatorname{last}(A),\operatorname{first}(B))$.

\begin{lemma}[Cut-pair extraction]
The union of these pairs over reachable binary rules is exactly $E(W)$.
It has at most as many entries as there are reachable binary rules and
can be collected without expanding $W$.
\end{lemma}
\begin{proof}
A rule's cut contributes an actual adjacent pair whenever that rule occurs
in the derivation tree. Conversely, any adjacent positions in the expanded
word have a lowest common ancestor in that tree. They lie on opposite sides
of its cut, so their pair is contributed by the corresponding grammar rule.
Repeated occurrences of a rule contribute the same pair and need not be
visited separately.
\end{proof}

Let $C$ be the shared signed alphabet and let $E=E(X)\cap E(Y)$ for the
two query roots. A common substring is a word over $C$ whose every adjacent
pair lies in $E$. It is a walk in the directed graph of allowed transitions,
although the implementation only needs a set of pairs and three lengths
per grammar rule. Signs and order are retained; inverse words transform
an original pair $(u,v)$ into $(-v,-u)$.

\subsection{Longest runs respecting the common transitions}

As before, summarize each grammar word by its longest valid prefix $p$,
suffix $q$ and contiguous run $r$. A terminal in $C$ has all three values
one, and any other terminal has all three zero. At a binary rule $W=UV$,
let $\delta=1$ when
$(\operatorname{last}(U),\operatorname{first}(V))\in E$, and zero otherwise.
The recurrence is
\[
\begin{aligned}
 p(W)&=p(U)+\delta\,\mathbf1_{p(U)=|U|}\,p(V),\\
 q(W)&=q(V)+\delta\,\mathbf1_{q(V)=|V|}\,q(U),\\
 r(W)&=\max\{r(U),r(V),\delta\,(q(U)+p(V))\}.
\end{aligned}
\]
A run crossing the cut must have an allowed transition there. If either
side contributes no valid letters, the child maxima already cover the
remaining run. Induction therefore gives exact run lengths, and
\[
 \operatorname{LCS}(U,V)\le\min(r(U),r(V))
\]
for any pair of reachable words in the query. The summaries use the shared
root constraints; they remain upper bounds for every smaller rule pair.

The implementation first uses prefix and alphabet-run checks. Only when
those fail to finish does it collect directed pairs. A side whose entire
pair set is shared needs no new transition-summary pass; its existing
alphabet bound remains valid. Otherwise the stronger summaries replace
its bounds. The existing incremental witness loop then uses the improved
global and local bounds, returning only after an exact witness reaches
one of them. This preserves the complete fallback and its polynomial
local bound. The extra extraction and summary work is linear in reachable
rules in the usual arithmetic/dictionary model; integer and dictionary
costs remain polynomial in the encoded input size.

For $X=(ab)^Nac$ and $Y=(ab)^Nca$, the shared pairs are $ab$ and $ba$.
The maximum valid run has length $2N+1$ in $X$ and $2N$ in $Y$.
The already checked prefix of length $2N$ therefore proves the exact
answer without any occurrence table. There is also an interior-match
family $c(ab)^Nb,b(ab)^Nc$. Cutting both grammars halfway through their
long periodic blocks lets the first aligned-cut extension find $2N$
letters. The transition bound is $2N$, so no further extensions are needed.
On the specified binary-power parses, these tests use polynomial work in
$\log N$ and no expansion.

\subsection{Filtering whole donors before matching}

An occurrence of a donor spelling $D$ inside a cyclic target $R$ requires
\[
 E(D)\subseteq E(R)\cup\{(\operatorname{last}(R),
                              \operatorname{first}(R))\}.
\]
The target's closing pair is essential: a genuine match may cross its seam.
For an inverse donor, transform every internal source pair as described
above. The test is only necessary; passing it still requires the exact
compressed occurrence query. A one-letter donor has no internal pairs
and is not rejected by this condition.

After the existing length and absolute-count filters, the implementation
collects each needed root's pairs lazily and caches them for that call.
It rejects incompatible orientations before allocating an inverse donor.
The total cached entries are limited by the arena's node allowance;
exhaustion raises the existing resource exception and remains inconclusive.
All traversals and candidate checks consume the shared work allowance.

For the cyclic relators
\[
 D=(ab)^Nac,\qquad R=(ab)^Nca,
\]
the absolute counts agree. However $ac$ is absent from the target's cyclic
pair set, and $bc$ is absent from the donor's cyclic pair set. The reversed
negative pairs cannot occur in either positive word. Thus every whole-donor
orientation is rejected without an occurrence table or inverse construction.
The complete cyclic stage then finds the common prefix $(ab)^N$ using
transition bounds on the doubled words. The ordinary relator move replaces
$R$ by $c^{-1}a^{-1}ca$, preserving the donor and normal closure. Total
expanded relator length falls from $4N+4$ to $2N+6$.

This measures a local search-and-application operation on supplied relators.
It does not establish that the optional group search reaches this state
before its other moves, nor does it certify an unknot from this presentation.
The public recognizer still reconstructs and independently checks a complete
trace from the original planar diagram.

\subsection{A control that the new condition does not resolve}

Let $T=\mathtt{aabbccacbabc}$. Its internal pairs include all nine ordered
pairs over $\{a,b,c\}$. The words
\[
 X=(ab)^N cT,\qquad Y=(ab)^N aT
\]
therefore have identical alphabets and identical pair sets. Their exact
LCS is $2N$ for $N\ge16$, but neither cheap bound proves this.
To see the exact value, the common prefix gives the lower bound. A common
substring longer than $2N$ would have to start within the first thirteen
positions of each word, since each tail has length thirteen. Unequal
starting parities mismatch immediately. With equal parity, if $X$ reaches
its tail first, its $c$ disagrees with the other alternating prefix. If
$Y$ reaches its tail first, its first $a$ can agree but its next $a$
disagrees with the following $b$. Equal starts disagree at $c$ versus $a$.
None yields a match longer than $2N$.

The complete fallback is still required on this control. Its retained
resource limits demonstrate that a stronger necessary condition does not
solve general compressed substring search by itself. More importantly,
none of these local bounds proves a bound on total relator moves, grammar
growth or recognition completeness. The general subexponential target
remains unresolved.

\subsection{Validation and controlled measurements}

Four new tests compare directed pair sets and transition-run summaries
with every reachable literal word in 400 random signed grammars. They check
same-alphabet queries through $N=2^{500}$ with occurrence tables forbidden,
an interior witness with later overlaps forbidden, an equal-pair control
requiring the exact fallback, cyclic-seam matches, the sign/order change
under inversion, impossible orientations with matching and inverse allocation
forbidden, an independent cache limit, exhausted work, and the complete
local partial-move operation without expansion. The integrated suite passes
686 tests with Regina 7.4.1 in 155.970 seconds; the fourteen-test targeted
suite passes in 20.006 seconds. A fifth new test forbids adjacency
traversal when uniform metadata already determines the pair set, including
one-letter and inverse-only matches. After that shortcut, the final
integrated suite passes 687 tests in 162.201 seconds.

\begin{center}\small
\begin{tabular}{@{}lrrr@{}}
\toprule
Scope & Before & A/A & Transitions \\
\midrule
Small 14, summed medians & 197.913 & 196.984 & 199.021 \\
Gordian, full query & 1947.100 & 1951.987 & 1941.596 \\
\bottomrule
\end{tabular}
\end{center}
Median milliseconds; all full knot queries complete.


All 225 full knot queries complete with identical certificate digests, and
none reaches the changed fallback. The small-fourteen sum of medians is
197.913 milliseconds before, 196.984 for A/A and 199.021 after. Gordian
medians are 1.947, 1.952 and 1.942 seconds. These data check existing paths;
they do not demonstrate a knot-corpus speedup.

\begin{center}\small
\begin{tabular}{@{}llrrr@{}}
\toprule
LCS family & $h$ & Before & A/A & Transitions \\
\midrule
Prefix & 8 & 1.089 & 1.099 & 0.122 \\
Prefix & 32 & 8.395 & 8.376 & 0.283 \\
Prefix & 500 & limit & limit & 3.612 \\
Interior & 8 & 85.855 & 87.077 & 0.197 \\
Interior & 32 & limit & limit & 0.463 \\
Interior & 500 & limit & limit & 5.120 \\
Equal bigrams & 8 & 7.066 & 7.085 & 7.235 \\
Equal bigrams & 32 & 71.284 & 71.929 & 71.226 \\
Equal bigrams & 500 & limit & limit & limit \\
\bottomrule
\end{tabular}
\end{center}
Median milliseconds including construction; limits are censored.
\begin{center}\footnotesize
\begin{tabular}{@{}llrrrrr@{}}
\toprule
Operation & $h$ & Before & A/A & Matcher & Pruning & Both \\
\midrule
Partial & 8 & 176.782 & 174.430 & 12.035 & 165.455 & 0.356 \\
Partial & 32 & limit & limit & 285.692 & limit & 0.762 \\
Partial & 500 & limit & limit & limit & limit & 7.078 \\
Quotient & 8 & 0.330 & 0.325 & --- & --- & 0.363 \\
Quotient & 32 & 0.909 & 0.930 & --- & --- & 1.025 \\
Quotient & 500 & 14.130 & 14.105 & --- & --- & 15.340 \\
\bottomrule
\end{tabular}
\end{center}
Median milliseconds including construction. Both denotes the initial filter before the uniform-word shortcut. Partial operations do not produce knot verdicts.


The primary audit contains 135 substring queries, 75 local partial-move
operations and 45 successful abstract quotient searches. At $h=500$, the
prefix family finishes in 3.612 milliseconds with 508 nodes and 6576 work
units; the interior family takes 5.120 milliseconds with 1014 nodes and
9596 work units. Both archived queries exhaust the two-million-work cap.
On the equal-bigram control the new stage adds little charged work but
cannot improve the bound: all arms still exhaust the cap at $h=500$.

For the full partial-move operation, at $h=8$ the baseline takes 176.782
milliseconds, the new matcher alone 12.035, pruning alone 165.455 and
both 0.356. At $h=500$, only both changes complete: 7.078 milliseconds,
514 nodes and 15775 work units. Both isolated components exhaust the cap.
This demonstrates why progress in a fallback must be accompanied by a
bounded route through the preceding stages on this family.

The successful quotient control exposes a regression in the initial
filter. At $h=500$, its median rises from 14.130 milliseconds (14.105 A/A)
to 15.340, and work increases from 27624 to 32135. The pair sets of these
uniform words were being recomputed by grammar traversal even though the
arena already stores their exact signed letter. The final implementation
uses that metadata: a one-letter word has no internal pair, and a longer
uniform word with letter $a$ has exactly the pair $(a,a)$. This takes
constant work and avoids the redundant traversal. Mixed words retain the
normal cut-based extraction, and cache accounting is unchanged.

\begin{center}\small
\begin{tabular}{@{}llrrrr@{}}
\toprule
Operation & $h$ & Before & A/A & Filter & Shortcut \\
\midrule
Quotient & 8 & 0.337 & 0.336 & 0.365 & 0.331 \\
Quotient & 32 & 0.950 & 0.952 & 1.062 & 0.969 \\
Quotient & 500 & 13.601 & 13.549 & 15.112 & 13.545 \\
Partial & 500 & --- & 7.329 & 7.445 & 7.346 \\
\bottomrule
\end{tabular}
\end{center}
Median milliseconds over 21 measured rounds. A/A repeats Before for Quotient and Filter for Partial. All operations complete with identical traces.


At $h=500$, the follow-up medians are 13.601 milliseconds before, 13.549
for its A/A control, 15.112 for the initial filter and 13.545 for the final
shortcut. The extra charged work falls from 4511 units to six; the node
count remains 1003. The nonuniform partial-operation medians are 7.445
milliseconds for the initial filter, 7.329 for its A/A control and 7.346
for the final code, all with the same trace, 514 nodes and 15775 work units.
Thus reuse of exact uniform metadata removes the measured quotient
regression without changing the mixed-word filtering behavior.

The follow-up compares the original donor helper, the initial filtered
helper archived at \code{3e9afde43}, and the final uniform shortcut in the
same current kernel and search harness. It uses 21 measured rounds after
one warmup, with shuffled arm order and seed 3011. There are 252 quotient
measurements and 63 partial-operation controls. The partial controls repeat
the initial filter as A/A and verify that the shortcut leaves nonuniform
traces unchanged. Every follow-up operation completes, and traces agree
across its arms. Raw follow-up samples are in
\path{results/donor_pair_shortcut_20261008.json}.


The baseline uses the matcher class and whole-donor helper archived at
\code{87e9585ec}, inside the common current word arena, recognition and
replay harness. A/A repeats those components. Local partial-move tests
also isolate the new matcher with the old whole-donor stage and the new
whole-donor filter with the old matcher. A successful pure-power quotient
family is retained to expose any overhead added to an already successful
whole-donor path.

Every case has one excluded warmup and five measured rounds with shuffled
arm order, using seed 3010. Full knot timings include fresh PD construction,
recognition and independent certificate replay; certificate digests are
checked outside the timer. Kernel and operation timings include grammar
construction. Resource-limited runs are censored, not negative answers
or completed-time denominators. Whole recognition allows 18 seconds,
with 15 seconds and twenty million work units in its group stage. Local
queries allow two million work units and 100000 nodes/table cells.
Raw results are in \path{results/lcs_transitions_20261008.json}, with
logs in \path{synthesis/data/lcs-transitions-*}. There was no new incoming
report, and the suggested \code{openai/math} head remained at the previously
recorded \code{fd4aeeb2ee}.
